\documentclass[11pt]{article}
\usepackage[margin=1in]{geometry}
\usepackage{amsmath,amssymb}
\usepackage{graphicx}
\usepackage[T1]{fontenc}
\usepackage[utf8]{inputenc}

\title{Part I: Early-Stage Convergence of Adam and SGD}
\author{}
\date{}

\begin{document}
\maketitle

\subsection{Part I: Does Adam Converge Faster Than SGD in Early-Stage Training?}

\paragraph{Motivation.}
Early-stage convergence is worth isolating because optimization trajectories are shaped near initialization, where different update rules can move the same network toward different regions of parameter space before later-stage behavior becomes visible. Accordingly, this part asks whether Adam reduces training loss faster than SGD at the beginning of training, whether it reaches the same loss thresholds in fewer updates, and whether the same pattern persists across architectures. We focus only on training dynamics here; questions about generalization, sharpness, Hessian-based proxies, and related geometric properties are deferred to later parts of the project.

\paragraph{Optimization Background.}
Let $\theta_t$ denote the parameter vector at step $t$, and let $g_t=\nabla_\theta \ell(\theta_t;B_t)$ be the minibatch gradient. SGD uses the update
\[
\theta_{t+1}=\theta_t-\eta g_t,
\]
so all coordinates share the same global learning-rate scale. Adam instead keeps exponential moving averages of the first and second moments,
\begin{align*}
m_t &= \beta_1 m_{t-1} + (1-\beta_1) g_t, \\
v_t &= \beta_2 v_{t-1} + (1-\beta_2) g_t^{\odot 2},
\end{align*}
with bias-corrected forms $\hat m_t$ and $\hat v_t$, and updates by
\[
\theta_{t+1}=\theta_t-\alpha \frac{\hat m_t}{\sqrt{\hat v_t}+\varepsilon}.
\]
The key distinction is therefore per-coordinate adaptive scaling: Adam rescales coordinates according to recent gradient magnitudes, whereas SGD does not. When early gradients are highly uneven across coordinates or layers, this normalization can make descent more effective during the transient phase of training. We use this as a mechanism-based explanation rather than a theorem claiming that Adam must always be faster.

\paragraph{Experimental Setup.}
We trained on Fashion-MNIST using two architectures: a batch-normalized MLP and a small CNN. All main experiments used batch size 128, 30 epochs, and cross-entropy loss. The default optimizer settings were SGD with learning rate 0.01 and momentum 0.9, and Adam with learning rate 0.001 and betas $(0.9,0.999)$. Results were aggregated over five seeds $\{42,123,456,789,2024\}$. We tracked training loss over epochs and steps, and recorded the first time the loss crossed thresholds 2.0, 1.0, 0.5, and 0.2. For each seed, Adam and SGD were matched by initialization and data order, so the comparison isolates optimizer behavior rather than randomness in starting point or minibatch sequence.

\paragraph{Results.}
\textit{MLP.} The MLP provides the strongest evidence for an early-stage advantage of Adam. Adam has lower mean training loss at every epoch, beginning already at epoch 1 ($0.4413$ vs.\ $0.4969$), and also ends slightly lower ($0.0552 \pm 0.0021$ vs.\ $0.0603 \pm 0.0017$). The threshold comparison is more decisive: at loss 0.2, Adam reaches the target in $450.4$ steps on average, whereas SGD requires $840.0$ steps. Moreover, Adam is faster at all four thresholds on all $5/5$ seeds, so the MLP result is both large and stable.

\textit{CNN.} The small CNN shows the same qualitative direction but much weaker evidence. Adam remains slightly ahead on average, yet the final losses are nearly tied ($0.0250$ vs.\ $0.0252$), and the mean threshold improvement is only about 5\%--16\%. Several seed-level comparisons are ties, including four ties out of five at threshold 0.2. The CNN experiment therefore supports the same conclusion only mildly.

\textit{Hyperparameter caveat.} The learning-rate search shows that Adam is clearly faster than the default SGD baseline, but the gap narrows when SGD is tuned more carefully. In particular, SGD with learning rate 0.03 is substantially stronger than the default 0.01 setting and can match Adam at some early thresholds. The most defensible claim is therefore not that Adam is universally faster than SGD, but that Adam has a robust advantage over the default SGD setting and a weaker, tuning-sensitive advantage more generally.

\begin{figure}[t]
\centering
\includegraphics[width=0.9\linewidth]{outputs/exp_a_mlp/plots/loss_vs_epoch.png}

\small \textbf{(a)} MLP loss vs.\ epoch.

\vspace{0.8em}

\includegraphics[width=0.92\linewidth]{outputs/exp_a_mlp/plots/threshold_comparison.png}

\small \textbf{(b)} MLP threshold comparison.

\caption{Main evidence for Part I from the MLP experiment. Panel (a) shows that Adam remains below SGD throughout training, with the clearest separation in the early epochs. Panel (b) shows that Adam reaches every loss threshold in fewer steps, with the largest practical gap at threshold 0.2.}
\label{fig:part1_mlp_main}
\end{figure}

\begin{figure}[t]
\centering
\begin{minipage}[t]{0.47\linewidth}
\centering
\includegraphics[width=\linewidth]{outputs/exp_b_cnn/plots/early_stage_zoom.png}

\small \textbf{(a)} CNN early-stage zoom.
\end{minipage}\hfill
\begin{minipage}[t]{0.50\linewidth}
\centering
\includegraphics[width=\linewidth]{outputs/exp_c_lr_search/plots/lr_search.png}

\small \textbf{(b)} Learning-rate search.
\end{minipage}

\caption{Supplementary evidence for Part I. Panel (a) shows that the Adam advantage persists for the CNN but is much smaller than for the MLP. Panel (b) shows that a stronger SGD learning rate substantially narrows the gap relative to the default SGD setting, so the conclusion should be interpreted with tuning sensitivity in mind.}
\label{fig:part1_supplementary}
\end{figure}

\paragraph{Discussion and Bridge to Part II.}
Taken together, the experiments support an \emph{early-stage optimization advantage} for Adam rather than a universal superiority claim. That advantage is strong on the MLP, much weaker on the CNN, and clearly sensitive to learning-rate choice. This interpretation is consistent with the adaptive-scaling mechanism above: Adam appears most helpful when the initial optimization transient is harder to negotiate, and much less distinctive when the architecture is already easy to optimize under both methods.

This result motivates the next part of the project. If Adam reaches low training loss faster from the same initialization, the natural follow-up question is whether the two optimizers merely differ in speed or whether they are being guided toward different kinds of solutions. Part II therefore asks whether faster early convergence is accompanied by a different implicit bias.


\end{document}
